% DD-C-0001; DD-C-0017
\subsection*{A research program in distributed discovery}

The sixteen-box model is the atomic case. A broader research program asks which distinctions survive when evidence is private, sources are correlated, actions overlap, search is sequential, and rewards are strategic. Seven directions are immediate. First, private-information teams can coordinate signal-contingent roles before observing their clues, filling the missing cell in the institutional matrix. Second, information design can ask whether a more informative public disclosure concentrates equilibrium search under a fixed strategic protocol. Third, latent source networks can replace the single common cue and make provenance overlap, rather than nominal report count, the unit of evidence diversity. Fourth, sequential discovery can allocate parallel and adaptive capacity after observed failures. Fifth, overlapping actions and multiple targets replace modular box coverage with set-valued or submodular discovery. Sixth, discovery mechanisms can jointly design reporting, action differentiation, and marginal-contribution credit. Seventh, empirical audits can estimate action concentration, source concentration, duplicate search, recovery budgets, and benchmark-relative protocol loss.

These are questions, not results. The framework supplies common accounting---information frontiers, protocol loss, budgets, redundancy, and provenance---but each extension needs its own feasibility class, objective, identification assumptions, and evidence status. In particular, the top-$L$ rule is special to atomic modular coverage, and the upstream effective-channel count is special to its common-cue process.
